% Part: computability
% Chapter: computability-theory
% Section: def-functions-self-reference

\documentclass[../../../include/open-logic-section]{subfiles}

\begin{document}

\olfileid{cmp}{thy}{slf}
\olsection{使用自指定义函数}

通过函数自身来定义函数通常很有用。例如，给定可计算函数 $k$、$l$ 和~$m$，不动点引理告诉我们，存在一个部分可计算函数~$f$，使得对每个~$y$ 都满足以下方程：
\[
f(y) \simeq
\begin{cases}
k(y) & \text{若 $l(y) = 0$} \\
f(m(y)) & \text{否则。}
\end{cases}
\]
同样，更具体地说，$f$ 是通过令
\[
g(x,y) \simeq
\begin{cases}
k(y) & \text{若 $l(y) = 0$} \\
\cfind{x}(m(y)) & \text{否则}
\end{cases}
\]
而得到的，然后应用不动点引理找到一个指标~$e$，使得
$\cfind{e}(y) = g(e,y)$。

作为一个具体的例子，“最大公约数”函数 $\fn{gcd}(u,v)$ 可以用下式定义：
\[
\fn{gcd}(u,v) \simeq
\begin{cases}
v & \text{若 $u = 0$} \\
\fn{gcd}(\fn{mod}(v, u), u) & \text{否则}
\end{cases}
\]
其中 $\fn{mod}(v, u)$ 表示 $v$ 除以~$u$ 的余数。应用不动点引理可知 $\fn{gcd}$ 是部分可计算的。（事实上，这可以写成上述形式，令 $y$ 编码序对 $\tuple{u, v}$。）随后对~$u$ 进行归纳即可证明，$\fn{gcd}$ 实际上是全函数。

当然，我们可以构造出比刚才讨论的例子复杂得多的自指定义。大多数编程语言都以某种方式支持用函数自身来定义函数。注意，与用计算该函数的算法的\emph{指标}来定义函数相比，这稍显平淡，而不动点定理在最一般的意义上允许我们做的正是后者。

\end{document}